% Source: https://artofproblemsolving.com/community/c3h2668968_proving_problems_thread
% Stable identity: aops-topic-2668968-q19
Count the ordered selections of $k$ different elements in two ways.
\begin{itemize}
\item Directly: there are $n$ choices for the first element, $n-1$ for the second, and so on, giving $n(n-1)\cdots(n-k+1)=\dfrac{n!}{(n-k)!}$.
\item Via subsets: choose the $k$-element subset ($\dbinom nk$ ways) and then order it ($k!$ ways).
\end{itemize}
Hence $\dbinom nk\cdot k!=\dfrac{n!}{(n-k)!}$, which gives the formula. (For $k=0$ both sides equal $1$.) $\blacksquare$
